% ============================================================================
\section{Operating Principle and Topological Characteristics}
\label{sec:operating_principle}
% ============================================================================

This section formalizes the operating principle underlying SPB's two
coupled properties -- reduced semiconductor voltage stress and
converter--machine modularity -- and establishes the notation and
voltage/power relations used throughout the remainder of the paper.

\subsection{Basic Structure and Voltage Partitioning}

An $N$-cell SPB consists of $N$ conventional three-phase two-level
voltage-source inverter cells connected in series on the dc side, with
each bridge supplying one electrically separated three-phase winding
group of the same machine \cite{Norrga2013SPB,Jin2014Control,
Jin2017Modulation}. The machine therefore presents $N$ ac ports, while
the converter cells share a single series dc path. This structure is
the physical implementation of the voltage-partitioning principle
introduced in Section~\ref{sec:semiconductor_limitations}.

The dc-link voltage satisfies

\begin{equation}
V_{\mathrm{dc}}=\sum_{k=1}^{N}V_{\mathrm{dc},k},
\qquad
V_{\mathrm{dc},k}\approx\frac{V_{\mathrm{dc}}}{N}
\label{eq:spb_voltage_sum}
\end{equation}

under balanced operation. A basic implementation contains $6N$ active
switches and $N$ local dc-link capacitors. Increasing $N$ therefore
reduces the nominal cell voltage, and hence the required semiconductor
voltage class, but simultaneously increases the number of switches,
gate drivers, sensors, capacitors, and machine-winding terminals that
must be realized and managed. This tradeoff between semiconductor
voltage rating and component count is quantified further in
Section~\ref{sec:semiconductor_scaling} and is revisited from a
machine-design perspective below. Fig.~\ref{fig:spb_switch_schematic}
shows this structure at switch level for two representative cells,
making explicit the $6N$-switch bridge realization, the local
capacitors $C_k$, and the notation used throughout the remainder of
this paper.

\begin{figure}[!t]
\centering
\resizebox{\columnwidth}{!}{%
\begin{circuitikz}[font=\scriptsize, x=1mm, y=1mm, >=latex]
\ctikzset{bipoles/length=5mm}

% ---- dc source (3-cell battery-pack symbol) ----
\draw (-10,10.5) -- (-10,28) -- (-4.5,28) -- (0,28);
\draw (-10,-10.5) -- (-10,-28) -- (-4.5,-28) -- (0,-28);
\draw (-10,10.5) to[battery1] (-10,3.5) to[battery1] (-10,-3.5)
  to[battery1] (-10,-10.5);
\node[font=\tiny] at (-5.3,0) {$V_{\mathrm{dc}}$};
\draw[<->] (-19,-28) -- (-19,28);
\node[rotate=90] at (-21.3,0) {$V_{\mathrm{dc}}$};
\node[font=\tiny,anchor=south] at (-4.5,28.3) {$+$};
\node[font=\tiny,anchor=north] at (-4.5,-28.3) {$-$};

% ---- series dc backbone (elided cells 2..N-1) ----
\draw (0,16) -- (0,9);
\node[font=\scriptsize] at (0,0) {$\vdots$};
\draw (0,-9) -- (0,-16);

% ---- series dc current ----
\draw[->,thick] (-2.3,25) -- (-2.3,19);
\node[font=\tiny,anchor=east] at (-2.7,22) {$i_{\mathrm{dc}}$};

% ================= CELL 1 =================
\node[font=\tiny] at (21,31.8) {Cell 1};
\draw (0,28) -- (40,28);
\draw (0,16) -- (40,16);
\draw (6,28) to[C, l=$C_1$] (6,16);
\node[font=\tiny,align=center] at (21,12.5) {$V_{\mathrm{dc},1}\approx V_{\mathrm{dc}}/N$};

\draw (13,28) to[nos] (13,22);
\draw (13,22) to[nos] (13,16);
\draw (21,28) to[nos] (21,22);
\draw (21,22) to[nos] (21,16);
\draw (29,28) to[nos] (29,22);
\draw (29,22) to[nos] (29,16);

\draw (13,22) -- (15.5,22) node[right=1pt,font=\tiny] {$a_1$};
\draw (21,22) -- (23.5,22) node[right=1pt,font=\tiny] {$b_1$};
\draw (29,22) -- (31.5,22) node[right=1pt,font=\tiny] {$c_1$};

% ---- star (wye) winding group 1, local neutral n_1 ----
% three separate lanes (well-separated in y) carry each phase's
% winding coil, converging only in a short diagonal stub at the end
% so the three coils remain visually distinct rather than bunching.
\draw[orange!70!black] (15.5,22) -- (15.5,26) -- (40,26);
\draw[orange!70!black] (40,26) to[cute inductor] (46,26);
\draw[orange!70!black] (46,26) -- (52,22);
\draw[orange!70!black] (23.5,22) -- (23.5,18) -- (40,18);
\draw[orange!70!black] (40,18) to[cute inductor] (46,18);
\draw[orange!70!black] (46,18) -- (52,22);
\draw[orange!70!black] (31.5,22) -- (40,22);
\draw[orange!70!black] (40,22) to[cute inductor] (46,22);
\draw[orange!70!black] (46,22) -- (52,22);
\node[circle,fill=orange!70!black,inner sep=0.5pt] at (52,22) {};
\node[font=\tiny,anchor=west,orange!70!black] at (53.3,22) {$n_1$};

\draw[->,blue!55!black] (2,31) -- (2,28.3);
\node[font=\tiny,anchor=south,blue!55!black] at (2,31.2) {$p_{\mathrm{dc},1}$};

% ================= elided cells =================
\node[font=\tiny,align=center] at (21,0) {(cells $2,\dots,N\!-\!1$\\elided)};
\node[font=\scriptsize] at (50,0) {$\vdots$};

% ================= CELL N =================
\node[font=\tiny] at (21,-31.8) {Cell $N$};
\draw (0,-16) -- (40,-16);
\draw (0,-28) -- (40,-28);
\draw (6,-16) to[C, l=$C_N$] (6,-28);
\node[font=\tiny,align=center] at (21,-12.5) {$V_{\mathrm{dc},N}\approx V_{\mathrm{dc}}/N$};

\draw (13,-16) to[nos] (13,-22);
\draw (13,-22) to[nos] (13,-28);
\draw (21,-16) to[nos] (21,-22);
\draw (21,-22) to[nos] (21,-28);
\draw (29,-16) to[nos] (29,-22);
\draw (29,-22) to[nos] (29,-28);

\draw (13,-22) -- (15.5,-22) node[right=1pt,font=\tiny] {$a_N$};
\draw (21,-22) -- (23.5,-22) node[right=1pt,font=\tiny] {$b_N$};
\draw (29,-22) -- (31.5,-22) node[right=1pt,font=\tiny] {$c_N$};

% ---- star (wye) winding group N, local neutral n_N ----
% mirrored version of the Cell 1 lane layout above.
\draw[orange!70!black] (15.5,-22) -- (15.5,-26) -- (40,-26);
\draw[orange!70!black] (40,-26) to[cute inductor] (46,-26);
\draw[orange!70!black] (46,-26) -- (52,-22);
\draw[orange!70!black] (23.5,-22) -- (23.5,-18) -- (40,-18);
\draw[orange!70!black] (40,-18) to[cute inductor] (46,-18);
\draw[orange!70!black] (46,-18) -- (52,-22);
\draw[orange!70!black] (31.5,-22) -- (40,-22);
\draw[orange!70!black] (40,-22) to[cute inductor] (46,-22);
\draw[orange!70!black] (46,-22) -- (52,-22);
\node[circle,fill=orange!70!black,inner sep=0.5pt] at (52,-22) {};
\node[font=\tiny,anchor=west,orange!70!black] at (53.3,-22) {$n_N$};

\draw[->,blue!55!black] (2,-13) -- (2,-16.3);
\node[font=\tiny,anchor=south,blue!55!black] at (2,-12.8) {$p_{\mathrm{dc},N}$};

\end{circuitikz}}
\caption{Switch-level realization of the SPB dc--ac stage for two
representative cells (Cell~1 and Cell~$N$; cells $2,\dots,N-1$
elided). Each cell is a conventional two-level three-phase bridge
(six switches per cell, $6N$ total) fed from its own local dc-link
capacitor $C_k$, with leg midpoints $a_k,b_k,c_k$ feeding a
star-connected winding group $k$ at its own electrically separate
neutral point $n_k$ (only $n_1$ and $n_N$ shown; the star points are
not interconnected). The common series current $i_{\mathrm{dc}}$
\eqref{eq:series_current} flows through every cell, giving each cell
the dc-side power $p_{\mathrm{dc},k}=V_{\mathrm{dc},k}i_{\mathrm{dc}}$
\eqref{eq:cell_dc_power}.}
\label{fig:spb_switch_schematic}
\end{figure}

\subsection{Series DC Connection and Power Flow}

The defining structural difference between SPB and parallel modular
drives is the series dc connection. Neglecting auxiliary and bypass
paths, the cell dc currents are constrained to be equal,

\begin{equation}
i_{\mathrm{dc},1}=i_{\mathrm{dc},2}=\cdots=i_{\mathrm{dc},N}
=i_{\mathrm{dc}},
\label{eq:series_current}
\end{equation}

so that the dc-side power entering cell $k$ is

\begin{equation}
p_{\mathrm{dc},k}=V_{\mathrm{dc},k}i_{\mathrm{dc}}.
\label{eq:cell_dc_power}
\end{equation}

The total machine power, by contrast, is the sum of the powers
delivered to the individual winding groups,

\begin{equation}
P_{\mathrm{ac}}=\sum_{k=1}^{N}P_k.
\label{eq:power_sum}
\end{equation}

Under symmetrical operation $P_k\approx P_{\mathrm{ac}}/N$, but equal
power sharing is not a fundamental requirement of the topology. The
multi-three-phase machine provides additional current and power
degrees of freedom that allow power to be redistributed among winding
groups while the required total torque is maintained
\cite{Che2014SeriesConnected,Verkroost2024PowerDistribution}. This
redistribution capability is the physical basis for both dc-link
voltage balancing and several of the post-fault operating strategies
discussed in later sections.

The combination of a common dc current with independently
controllable ac-side powers also explains why the local capacitor
voltages must be treated as internal dynamic states rather than fixed
design parameters. If one winding group processes more power than
another, the corresponding capacitor energies can diverge even though
all cells carry the same dc current \cite{Nikouie2015Stability,
Nikouie2017Stability}. Managing this divergence is precisely the
dc-link stability and balancing problem identified in
Section~\ref{sec:evolution} as the topology's first-order control
challenge.

\subsection{SPB as an Input-Series, Multi-Three-Phase-Output Drive}

These relations place SPB in a distinct category relative to the
conventional multilevel converters discussed in
Section~\ref{sec:introduction}, which partition dc-link voltage
internally while retaining a conventional three-phase machine. In
SPB, by contrast, each converter cell remains a standard two-level
bridge, and it is the machine itself that provides the multiple ac
ports required to transfer power out of the series dc stack. SPB can
therefore be characterized as an \emph{input-series,
multi-three-phase-output} drive, a description that makes explicit
why voltage partitioning in SPB cannot be separated from machine
design.

\subsection{Cell Count as a Converter--Machine Design Variable}

Cell count is consequently a converter--machine design variable rather
than a semiconductor variable alone. As established above, increasing
$N$ also reduces the voltage available to each winding group and
alters the number of winding sets, terminal count, electromagnetic
coupling, current-control degrees of freedom, and available
reconfiguration states. Bringezu and
Biela demonstrated this interaction directly by comparing SPB systems
based on 3-, 6-, 9-, and 12-phase PMSMs \cite{Bringezu2020SPB}.

A related and often overlooked point is that machine phase current
does not generally scale as $1/N$. The required phase current instead
depends on the number of turns, flux linkage, winding connection,
torque constant, and machine geometry. A many-cell architecture may
therefore reduce the semiconductor voltage class while still requiring
high-current devices or parallel chips at the device level. This
coupling between cell count, voltage class, and current rating is
particularly important when the comparison is extended from
high-voltage SiC implementations toward lower-voltage GaN
implementations, as considered in Section~\ref{sec:semiconductor_limitations}.
